% =============================================================================
% 047 / 068
% Status: raw
% =============================================================================
% Notes:
%
% Source question:
% 先不論吳語未來如何，
% 請問你怎麼看 受到官話思維影響的諸語言，
% 未來會如何發展？官話本身未來又會如何？
% 使用這些語言的人，在可預期的未來，
% 將因官話思維深受其害、而不自知？
% =============================================================================

\chapter[先不論吳語未來如何， 請問你怎麼看 受到官話思維影響的諸語言， 未來會如何發展？官話本身未來又會如何？ 使用這些語言的人，在可預期的未來， 將因官話思維深受其害、而不自知？]{先不論吳語未來如何， \\[0.35em] 請問你怎麼看 受到官話思維影響的諸語言， \\[0.35em] 未來會如何發展？官話本身未來又會如何？ \\[0.35em] 使用這些語言的人，在可預期的未來， \\[0.35em] 將因官話思維深受其害、而不自知？}
\qaanswer{
首先官話和普通話不是同一個概念，儘管英語對於兩者的解釋都是Mandarin，它包含了官話與普通話。普通話是官話的2.0版，融入了更多英語語法的建構成分(儘管這種吸收和建構非常失敗)。Mandarin未來能發展到如何我無法預言，但我相信以這種語言為母語的人群在思維邏輯上以及思辨能力肯定是不如屈折語和黏著語母語的人群的。妳不可能讓一個對時態、語態、事態、格、屬、性毫無概念的人能在根本的思維層面上明白和運用這些，而這些又和邏輯思維、思辨力息息相關。簡單來說Mandarin到目前為止都沒有很好的整理出一套可以自圓其說的語法。所以以Mandarin為母語的文明中出現思維精神分裂的縫合怪非常的司空見慣。不借助其他邏輯思辨更強的語言作為思維工具，單憑Mandarin甚至連學習高等數學都會十分吃力，成才率也更低。從這個方面來說，大眾爭論不休的是否廢除漢字的討論其實倒是很次要的問題。東亞諸國書面語離開漢字是非常不現實的，但用全漢字作為書面語表述工具是會阻礙語法建構和思維發展的。
}

